\chapter{Pulled-tight sequences and the limits of stationarity}\label{sec:pulled-tight-limits}
The preceding deformation now acts on whole minimizing sequences. The compactness argument proves stationarity without assuming rectifiability of an arbitrary weak varifold limit. Almost-minimizing selection and replacement regularity remain further steps.

\begin{definition}[A pulled-tight minimizing sequence]\label{def:pulled-tight-sequence}
For a fixed flat-homotopy class $\Pi$ of $\mathbf F$-continuous maps on a finite simplicial complex, a minimizing sequence is called pulled-tight if every member of its critical set is stationary. This is \cite{MNIndex}, Definition 2.9. The terminology records stationarity, not smoothness, multiplicity one or index control.
\end{definition}

\begin{theorem}[Pull-tight preserves the width and decreases the critical set]\label{thm:pulled-tight-minimizing-sequence}
Let $M^{n+1}$ be smooth, closed and connected, $n\geq1$, and let $\Phi_i\in\Pi$ be a minimizing sequence in the mod-two boundary-cycle setting of Definition \ref{def:homotopy-width}, with finite width $L$. There is a minimizing sequence $\Psi_i\in\Pi$ such that
\[
 \mathcal C(\Psi_i)\subset\mathcal C(\Phi_i)\cap\{V:\delta V=0\}.
\]
Moreover $\Psi_i$ may be chosen uniformly as close to $\Phi_i$ in $\mathbf F$ as prescribed. This is the continuous pull-tight conclusion of \cite{MNIndex}, Section 2.8; the discrete version discussed in \cite{LiAP}, Section 4.2, additionally requires interpolation.
\end{theorem}
\begin{proof}
Choose $C>\sup_{i,x}\mathbf M(\Phi_i(x))$, finite because the maximal masses converge. Use one fixed homotopy from Theorem \ref{thm:continuous-pulltight-homotopy} and set $\Psi_i(x)=H(1,\Phi_i(x))$. The composition is $\mathbf F$-continuous and the homotopy $H(s,\Phi_i(x))$ is flat-continuous, so $\Psi_i\in\Pi$. By the definition of the width and mass decrease,
\[
 L\leq\max_x\mathbf M(\Psi_i(x))\leq\max_x\mathbf M(\Phi_i(x))\longrightarrow L.
\]
Let $W$ be critical for the new sequence, represented by $|\Psi_{i_j}(x_j)|\to W$ with masses tending to $L$. Pass to a subsequence with $V_j=|\Phi_{i_j}(x_j)|\to V$ in $K_C$. The old masses lie between the new masses and the old maximal masses, so they also tend to $L$. Thus $V\in\mathcal C(\Phi_i)$ and the mass losses tend to zero. Corollary \ref{cor:pulltight-uniform-loss} gives $V\in S_C$. Continuity of the varifold deformation then gives
\[
 W=\lim_j\mathscr H(1,V_j)=\mathscr H(1,V)=V.
\]
This proves both stationarity and inclusion in the original critical set. The prescribed closeness follows from the same homotopy, without letting its positive defect vary with $i$.
\end{proof}

\begin{corollary}[Positive width produces a nonzero stationary varifold]\label{cor:positive-width-stationary-varifold}
Under the hypotheses of Theorem \ref{thm:pulled-tight-minimizing-sequence}, if $L>0$, there is a nonzero stationary $n$-varifold on $M$ of mass $L$ in the critical set of the original minimizing sequence.
\end{corollary}
\begin{proof}
The pulled-tight sequence has a nonempty critical set by Lemma \ref{lem:compact-critical-set}. Choose a member. Theorem \ref{thm:pulled-tight-minimizing-sequence} makes it stationary and a member of the original critical set; Lemma \ref{lem:compact-critical-set} gives mass $L>0$. This proof yields neither integrality nor regularity. Those stronger conclusions require additional variational selection, as in the external Theorem \ref{thm:ap-realization}.
\end{proof}

\begin{proposition}[Near-maximal slices of a pulled-tight sequence approach stationarity]\label{prop:pulledtight-near-maximal-stationary}
For a pulled-tight minimizing sequence with finite width $L$, every $\eta>0$ admits $\delta>0$ and $i_0$ such that
\[
 i\geq i_0,\quad \mathbf M(\Phi_i(x))\geq L-\delta
 \quad\Longrightarrow\quad d(|\Phi_i(x)|,S_C)<\eta,
\]
where $C$ bounds all masses. No rate in terms of $L-\mathbf M(\Phi_i(x))$ is asserted.
\end{proposition}
\begin{proof}
The open $\eta$-neighborhood of $S_C$ contains the critical set because the sequence is pulled-tight. Apply Lemma \ref{lem:near-critical-slices} to this neighborhood. The compactness argument supplies qualitative thresholds, with no quantitative lower bound on the auxiliary defect from first variation alone.
\end{proof}

\begin{example}[A diffuse stationary varifold on a compact manifold]\label{ex:diffuse-stationary-flat-torus}
For $n\geq1$ on the unit flat torus $M=\mathbb R^{n+1}/\mathbb Z^{n+1}$, let $P$ be the parallel plane spanned by the first $n$ coordinate fields. The varifold defined by
\[
 \int f\,dV=\int_M f(x,P)\,d\operatorname{Vol}_{n+1}(x)
\]
is stationary, has positive finite mass, and is not $n$-rectifiable.
\end{example}
\begin{proof}
For a smooth periodic vector field $X$, its first variation is $\int_M\sum_{i=1}^n\partial_iX_i\,dx=0$ by integration over each periodic coordinate. Its weight is ambient volume. Any countably $n$-rectifiable set has zero $(n+1)$-dimensional volume: each bounded piece of a Lipschitz image of a subset of $\mathbb R^n$ has zero ambient volume by the elementary dimension estimate, and countable unions retain this property. Thus ambient volume cannot be carried by such a set. This example shows why the compact mass cap in the deformation argument cannot silently be treated as a compact space of integral varifolds.
\end{proof}

\begin{example}[Stationary integral varifolds can have crossings]\label{ex:crossing-stationary-tori}
On the unit flat $(n+1)$-torus, $n\geq1$, the varifold sum of the coordinate tori $\{x_1=0\}$ and $\{x_2=0\}$ is integral and stationary. Its support is not a smooth embedded hypersurface near their intersection.
\end{example}
\begin{proof}
Both tori are closed and totally geodesic, hence stationary, and first variation is linear in the varifold. Their intersection has dimension $n-1$ and zero $n$-dimensional measure, so adding them gives an integral varifold with multiplicity one away from the intersection. At an intersection point its tangent cone is the sum of two distinct transverse $n$-planes. The support therefore cannot be locally a single smooth hypersurface, nor a disjoint collection of embedded hypersurfaces. Integrality and stationarity alone cannot replace the almost-minimizing and replacement arguments in Remark \ref{rem:pitts-proof-route}.
\end{proof}
